% ID: SEL-D0070-OB-E01
% ESTADO: cotejada. Segunda pasada visual y borrador compilado revisados.
% Fuente: Selectividad/sel_2009_matematicas_aplicadas/mat apli CC SS04.pdf
% SHA-256: bdea2fa323085e0db4eacff922ba9f403f58ef4e6bb4f2b569a0b857554c9ef2
% Página: 2; opción B; ejercicio original 1.

\textbf{(3 puntos)} Un agricultor posee 10 hectáreas (ha.) y decide dedicarlas al cultivo de cereales y hortalizas. Por las limitaciones de agua no puede destinar más de 5 ha. a hortalizas. El cultivo de cereales tiene un coste de 1000 euros/ha. y el de hortalizas de 3000 euros/ha., no pudiendo superar el coste total la cantidad de 16000 euros. El beneficio neto por ha. de cereales asciende a 2000 euros y el de hortalizas a 8000 euros. Halle la distribución de cultivos que maximiza el beneficio y calcule dicho máximo.
